\section{Overview of the \gai Framework}
\label{section:graph-based-ai-framework}

    
    \begin{figure}[h!]
        \begin{center}
            \includesvg[width=\columnwidth]{SVGs/gai_benchmark_high_level.svg}
            \caption{The high-level view of \gai benchmarking framework}
            \label{figure:gai-framework-overview}
        \end{center}
    \end{figure}



The \gai Benchmark Framework shown in Figure~\ref{Figure:gai-framework-overview} is an Input/Output benchmark that evaluates the performance of a system under test by providing problems, recording the time to solution, and verifying the correctness of an answer. 
The ideal system under test achieves a perfect effectiveness score in minimal time and cost, but all three dimensions will vary in the systems assessed. 

The benchmark is designed so that the system under test is completely independent of the benchmark itself. 
Using the Input/Output paradigm only requires the \ac{SUT} to interact with an \ac{API} (the observation apparatus). It so places no constraints on how a \ac{SUT} conFigures its hardware, software stack, or specific \gai implementation to solve the problem. 

The benchmark definition of the \gai framework is comprised of independent benchmark problem modules. 
Each module contains a collection of datasets designed to stress the efficiency and effectiveness of \ac{SUT} and uses \textit{problem-centric} task definitions to ensure problems are grounded in the real world and to encourage innovative solutions by removing the constraints on kernel-level tasks most existing benchmarks enforce. 

The independence of the benchmark specification from the \ac{SUT} means that an analyst, engineer, or researcher using the benchmark can quickly modify their \ac{SUT} and validate performance changes on a known benchmark. 
The careful design of the component benchmark problems to achieve the `grey-box' analysis of a \gai \ac{SUT} requires a formal methodology for transparency, reproducibility, and verifiability. 
